\documentclass{article}
\usepackage[utf8]{inputenc}
\usepackage[russian]{babel}
\usepackage{amsmath}
\usepackage{amsfonts}

\title{ADR-0002: Software Transactional Memory (STM) для управления конкурентным состоянием}
\author{Оперативная группа проекта SMP}
\date{Июль 2026}

\begin{document}

\maketitle

\section*{Статус}
\textbf{ACCEPTED} (Принято к исполнению)

\section*{Контекст}
Нашему L7-балансировщику \texttt{stm-markov-proxy} жизненно необходимо вести учет времени и распределять входящие сетевые запросы по пулу бэкендов в циклическом режиме (Round-Robin). Это означает, что тысячи независимых грин-потоков будут одновременно ломиться к одной и той же переменной-счетчику, пытаться прочитать текущий индекс, увеличить его на единицу и записать обратно.

Если использовать классические блокировки из мира Unix (мьютексы или семафоры), мы неизбежно столкнемся с двумя проблемами. Во-первых, блокировки не композируются — их тяжело объединять в сложные цепочки без риска намертво заклинить систему (Deadlock). Во-вторых, потоки будут тратить процессорные такты на ожидания в очереди за «замком», что убьет производительность под обстрелом утилиты стресс-тестирования \texttt{wrk}.

\section*{Решение}
Вместо пессимистичных блокировок утверждается использование \textbf{Software Transactional Memory (STM)}. Вся работа с динамическим счетчиком переносится в стерильную рантайм-лабораторию GHC:
\begin{enumerate}
    \item Глобальное изменяемое состояние заворачивается в транзакционную переменную \texttt{TVar ProxyState} и помещается внутрь общего регистра \texttt{Env}.
    \item Любое изменение счетчика происходит внутри изолированной монады \texttt{STM} через механизм TLog ( Transaction Log). Поток делает копию данных в свой локальный «черновик» и крутит логику там, не выставляя никаких замков физической памяти и никому не мешая.
    \item Мостом в реальный мир эффектов ввода-вывода становится атомарный комбинатор:
    \[ \text{atomically} :: STM\ a \to IO\ a \]
    Рантайм Haskell проверяет черновик на финише через процессорную инструкцию \texttt{CAS} (Compare-And-Swap). Если за время вычислений другой поток успел обновить счетчик раньше, наш черновик тихо уничтожается (\text{Rollback}), и вся транзакция детерминированно запускается заново с первой строки (\text{Retry}), гарантируя кристальную чистоту внешнего мира.
\end{enumerate}

\section*{Последствия}
\begin{itemize}
    \item \textbf{Плюсы:} Полная математическая защита от дедлоков под высокой нагрузкой. Жесткая изоляция эффектов на уровне системы типов GHC — компилятор физически не пропустит сетевой или консольный \texttt{IO}-код внутрь транзакции \texttt{STM}. Высокая скорость за счет работы с персистентными указателями без глубокого копирования объектов в памяти.
    \item \textbf{Минусы:} Риск высокой частоты перезапусков транзакций (\text{Transaction Contention}) при экстремальном конкурентном доступе на запись к одной переменной, что требует обязательного аудита под нагрузкой.
\end{itemize}

\end{document}
